Autoregressive language models train on all tokens in parallel, yet must generate them one at a time.
Diffusion language models remove this limit by denoising many tokens in each forward pass \citep{nie2026large, ye2025dream7bdiffusionlarge, team2026diffusiongemma}.
One promising area for research is diffusion language model samplers, as many advances in the continuous domain have been driven by improved sampler performance \citep{song2020denoising}.
The same is true for the language domain, where sampling order has been shown to have a dramatic effect on task performance \citep{kim2025train, ben2026accelerated, kim2026klass}.

Any diffusion language sampler makes two choices for each position in its denoising window.
The first is \emph{where} the next token comes from: the model's prediction, the token already at that position, or fresh noise.
This defines when a position within a denoising window has converged.
The second is \emph{when} the position is committed to the KV cache, after which it can no longer change.
These design choices give rise to two potential extremes: on one end, it is possible to commit tokens on a \textit{per position} basis, or as soon as a position has converged.
It is also possible to commit tokens \textit{at a block level}, or after the entire denoising window has converged.

In this work, we use the Entropy Bound sampler used in DiffusionGemma \citep{team2026diffusiongemma} to study these extremes and uncover the underlying tradeoffs.
Committing tokens on a per position level removes the chance for later revision, which is particularly useful when trading off performance for speed at higher entropy bounds.
When accepted tokens are frozen, macro accuracy falls from 73\% to 20\% as the entropy bound is relaxed, while the unfrozen sampler stays near 70\%.
Committing on a block level results in many redundant computations, as tokens converge in a largely left-to-right order.
Thus a converged prefix sits idle in the window, and by twenty denoising steps more than half of the slots, on average, are spent on tokens that are already final.

As either extreme results in losing sampler flexibility or tolerating wasted computation, we introduce the \textbf{R}olling \textbf{E}ntropy \textbf{S}ampler (RES) to strike the balance between the two.
Rather than committing tokens as soon as they have converged or waiting for the entire window, the sampler waits until a contiguous prefix of positions of length at least $s$ has converged.
Beyond avoiding redundant computations on converged tokens, this still allows for a smooth tradeoff between performance and speed via the EB threshold.

In summary, our contributions are as follows:
\begin{itemize}
    \item We write samplers as a routing rule and a commitment rule (Section~\ref{sec:motivation}), which separates accepting a token from committing it.
    \item We show that DiffusionGemma self-corrects without renoising, and that freezing accepted tokens destroys its speed-quality frontier (Section~\ref{sec:too-early}).
    \item We show that the EB sampler converges left to right, leaving over half of the window idle late in each block (Section~\ref{sec:too-late}).
    \item We introduce RES (Section~\ref{sec:method}). Across seven benchmarks it raises tokens per forward pass by 32.5\% at an average cost of 1.9 points of accuracy, and it extends the frontier beyond what relaxing the entropy bound alone reaches (Section~\ref{sec:experiments}).
\end{itemize}
